\chapter{OBJETIVOS}

\section{Alcance del proyecto}

Este proyecto abarca el diseño, desarrollo y validación de un sistema desacoplado de microservicios para la ingesta, procesamiento, y análisis operacional de datos de dispositivos IoT, utilizando como caso de uso la supervisión y análisis en tiempo real de la telemetría de consumo energético en edificios con medidores de electricidad, agua fría y agua caliente.

\section{Tareas a desarrollar}

Para cubrir el alcance descrito, se plantean las siguientes tareas:

\begin{itemize}
    \item Analizar el estado del arte en las tecnologías, estrategias y retos actuales en el procesamiento y análisis de datos en flujo.
    \item Analizar las herramientas de código abierto disponibles para el procesamiento y análisis de datos de dispositivos IoT.
    \item Diseñar e implementar un simulador de dispositivos IoT para la transmisión de datos de telemetría.
    \item Diseñar e implementar una arquitectura basada en contenedores para la orquestación y ejecución de los siguientes microservicios:
    \begin{itemize}
        \item Servicio para la ingesta de datos que acepte protocolos de telemetría estándar.
        \item Servicio de gobernanza y validación de esquemas de datos en flujo.
        \item Servicios de agregado y transformación que permitan calcular índices de eficiencia energética, efectuar análisis operativo e histórico y hacer detección de anomalías.
        \item Servicio de almacenamiento optimizado para series temporales.
        \item Servicio de almacenamiento optimizado para consumo analítico.
    \end{itemize}
    \item Diseñar e implementar servicios de visualización que permitan la exploración interactiva de los datos procesados.
\end{itemize}

\section{Metodología de desarrollo}
\label{chap:objetivos}

Las tareas del proyecto se desarrollarán como una serie de objetivos incrementales, acompañados de un indicador clave de rendimiento (KPI) que permite verificar su cumplimiento, y de las acciones necesarias para alcanzarlo. Los valores numéricos de cada KPI se han fijado a partir de rangos razonables de la práctica de ingeniería y de la literatura revisada en el estado del arte; se validarán y se recalibrarán de ser necesario, a partir de las pruebas realizadas.

\subsection{Objetivo 1: Garantizar la ingesta fiable de telemetría en tiempo real}

\textbf{KPI}: latencia inferior a 2 segundos para la ingesta del 95\,\% de los mensajes MQTT publicados, calculada desde la publicación del mensaje hasta su persistencia final, con una tasa de pérdida inferior al 0.1\,\%.

\textbf{Acciones}: instrumentar el simulador MQTT y desplegar Mosquitto como broker MQTT con un microservicio de bridge MQTT-Kafka que enviará los eventos ingeridos al tópico de Kafka correspondiente; definir el nivel de calidad de servicio (QoS) y la topología de tópicos; medir la latencia extremo a extremo.

\subsection{Objetivo 2: Garantizar la gobernanza del esquema de datos}

\textbf{KPI}: el 100\,\% de los eventos ingeridos se validan correctamente contra el esquema Avro registrado con 0\,\% de fallos de deserialización.

\textbf{Acciones}: definir el esquema Avro inicial de telemetría; configurar las reglas de compatibilidad de esquema en Apicurio; diseñar y ejecutar una prueba verificando que no se interrumpe el flujo ni se pierden mensajes.

\subsection{Objetivo 3: Procesar y enriquecer los datos en flujo con baja latencia}

\textbf{KPI}: cada micro-lote de Spark Structured Streaming se procesa en menos de 1 segundo sin acumular retraso entre lotes sucesivos.

\textbf{Acciones}: implementar el job de Spark Structured Streaming (DataFrame API) con ventanas de agregación temporal y \textit{watermarking} para el tratamiento de datos tardíos; medir la máxima duración de procesamiento cada micro-lote.

\subsection{Objetivo 4: Ofrecer visualización diferenciada operacional y analítica}

\textbf{KPI}: visualización mediante dashboards de Grafana y páginas de Power BI que muestren datos con una latencia menor o igual a 5 segundos cubriendo consumo energético, eficiencia operativa, detección de anomalías y análisis histórico.

\textbf{Acciones}: implementar los dashboards operacionales de Grafana sobre los datos de TimescaleDB; implementar las páginas analíticas de Power BI sobre los datos en PostgreSQL; definir y calcular los datos requeridos por los tipos de análisis; validar la latencia de actualización de las visualizaciones.

\subsection{Objetivo 5: Validar la resiliencia del sistema}

\textbf{KPI}: ante el fallo de un servicio individual, el sistema se recupera en menos de 60 segundos sin pérdida de datos.

\textbf{Acciones}: simular la caída de un servicio (por ejemplo, el broker de mensajería o una de las bases de datos) y medir el tiempo de recuperación y la tasa de pérdida de datos.

\section{Beneficios del proyecto}

El desarrollo de este proyecto aporta un trabajo académico evaluable que documenta la implementación completa de una arquitectura Kappa para IoT basada en herramientas de código abierto, cerrando parcialmente los vacíos identificados en el estado del arte respecto a la gobernanza de esquemas y la separación de sumideros operacional y analítico. 

Adicionalmente, el pipeline resultante queda disponible como sistema replicable sin dependencia de proveedores propietarios, lo que facilita su adopción como punto de partida en proyectos similares de monitorización energética.
